\documentclass[exercice]{thedocument}%
\startdatakeys
\source{}
\cycle{}
\competencesDuSocle{}
\connaissancesRequises{}
\competencesTravaillees{}
\enddatakeys
\begin{document}%
En utilisant l'inégalité triangulaire sur la figure ci-dessous, écrire six
inégalités différentes.
    \begin{center}
        \begin{tikzpicture}[scale=0.8]
            \coordinate(A)at(0,0);
            \coordinate(B)at(5,-2);
            \coordinate(C)at(2,-3);
            \coordinate(D)at(-1.5,-2.5);
            \draw node[right]at(B){\sf\strut B};
            \draw node[above]at(A){\sf\strut A};
            \draw node[below]at(C){\sf\strut C};
            \draw node[left]at(D){\sf\strut D};
            \draw(A)--(B)--(C)--(D)--cycle;
            \draw(A)--(C);
        \end{tikzpicture}
    \end{center}
\end{document}
